\subsubsection{Correlation on RGB Images}\label{sec:correlation-on-rgb-images}

Until now we assumed that an image was a single matrix:
\begin{equation*}
    I \in \mathbb{R}^{R \times C}
\end{equation*}
Where every pixel contains only one values (its intensity). However, real images are usually \textbf{RGB images}, so each pixel contains \textbf{three numbers}: Red (R), Green (G), and Blue (B) intensities. Therefore, an RGB image can be represented as a three-dimensional array:
\begin{equation*}
    I \in \mathbb{R}^{R \times C \times 3}
\end{equation*}
Instead of one image, we actually have three matrices, one for each color channel:
\begin{equation*}
    I_R \in \mathbb{R}^{R \times C}, \quad I_G \in \mathbb{R}^{R \times C}, \quad I_B \in \mathbb{R}^{R \times C}
\end{equation*}
Each channel stores the intensity of a single color.

\highspace
\textcolor{Green3}{\faIcon{question-circle} \textbf{Thus, does the correlation formula change for RGB images?}} Fortunately, \textbf{no}! The idea of correlation remains exactly the same. For grayscale images we had:
\begin{equation*}
    T\left(I\right)\left(r,c\right) = \sum_{u,v} w\left(u,v\right)I\left(r+u,c+v\right)
\end{equation*}
This computes a weighted sum over the neighborhood. For RGB images, every pixel has \textbf{three components}, so the filter must also have \textbf{three channels}. Instead of a single matrix $k \times k$, the filter becomes a three-dimensional array $k \times k \times 3$ because every position stores three weights, one for each color channel.

\highspace
\begin{flushleft}
    \textcolor{Green3}{\faIcon{book} \textbf{An RGB Filter}}
\end{flushleft}
A filter is no longer one matrix. It is actually three matrices, one for each color channel:
\begin{equation*}
    \mathbf{w} = \left\{ w_R, w_G, w_B \right\}, \quad w_R \in \mathbb{R}^{k \times k}, \quad w_G \in \mathbb{R}^{k \times k}, \quad w_B \in \mathbb{R}^{k \times k}
\end{equation*}
Together they form a three-dimensional array:
\begin{equation*}
    \mathbf{w} \in \mathbb{R}^{k \times k \times 3}
\end{equation*}
Now, the \textbf{correlation is computed channel-wise and summed}. Thus:
\begin{enumerate}
    \item Correlate the Red channel of the image with the Red channel of the filter.
    \item Correlate the Green channel of the image with the Green channel of the filter.
    \item Correlate the Blue channel of the image with the Blue channel of the filter.
    \item Sum the three results to obtain the final output.
\end{enumerate}
Mathematically, we \textbf{extend} the compact \textbf{correlation} formula (\autopageref{eq:linear-filter}) introducing an \textbf{additional summation over the color channels}:
\begin{equation}\label{eq:linear-filter-rgb}
    T\left(I\right)\left(r,c\right) = \sum_{i}\sum_{u,v} w_{i}\left(u,v\right) I_{i}\left(r+u,c+v,i\right)
\end{equation}
Where $i \in \left\{R,G,B\right\}$ is the color channel (e.g., $i=1$ for Red, $i=2$ for Green, and $i=3$ for Blue). Compared with the grayscale case, the only new operation is the outer sum over the three channels.

\highspace
\textcolor{Green3}{\faIcon{question-circle} \textbf{Why is this useful?}} Imagine we want to detect \emph{red objects}. The filter could learn \emph{red weights} are high, while \emph{green and blue weights} are zero. Then only regions rich in red produce a high response. Similarly, a filter detecting the sky might assign high weights to the blue channel and low weights to the red and green channels. This allows filters to learn \textbf{color-sensitive features}, not just shapes or edges.